\section{`We're a rich city with poor people': municipal strategies of new-build gentrification in Rotterdam and Glasgow}

\begin{description}
  \item[Authors] Doucet, B.; van Kempen, R.; van Weesep, J.
  \item[Year] 2011
  \item[Journal] Environment and Planning A, 43(6), 1438--1454
  \item[DOI] 10.1068/a43470
\end{description}

\subsection*{Domain Classification}

\begin{itemize}
  \item \textbf{Primary domain(s)}: Geography
  \item \textbf{Secondary domain(s)}: Regional Science (comparative institutional/political-economic analysis of urban governance regimes); tangential Urban Economics (urban competitiveness and housing-stock-mismatch discourse, without formal welfare or supply-elasticity modelling)
  \item \textbf{Keywords}: new-build gentrification, rent gap theory, urban entrepreneurialism, developer-led redevelopment, waterfront brownfield redevelopment, comparative case-study design, municipal policy discourse, socioeconomic upgrading, exclusionary displacement, elite semi-structured interviews
\end{itemize}

---

\subsection*{Research Question}

Why and how do municipal governments pursue new-build gentrification as an explicit urban policy goal, and how does this goal formation and its implementation vary across different political-economic (institutional) contexts? The paper addresses this through a comparative study of two flagship waterfront redevelopments -- the municipally led Kop van Zuid (Rotterdam) and the developer-led Glasgow Harbour (Glasgow) -- asking specifically what role new-build projects play in realizing this goal and why the rationale behind the strategy has received little direct empirical attention from the perspective of policymakers themselves.

\subsection*{Theoretical Framework}

The paper is situated within the gentrification literature originating with Glass's (1964) coinage of the term, tracing its extension from a description of individual household-level neighborhood upgrading toward a wider urban redevelopment strategy encompassing flagship and megaproject-scale ``new-build gentrification'' (Davidson and Lees, 2005). It explicitly engages the debate between Boddy (2007), who argues new-build development should not be classified as gentrification because it lacks direct displacement and an authentic historic aesthetic, and Davidson and Lees (2005), who treat new-build housing as one of the ``mutations'' of gentrification because it is driven by the closure of Neil Smith's (1979) \emph{rent gap}, produces indirect/exclusionary displacement, replicates a familiar gentrification aesthetic, and attracts a socioeconomically similar ``urban seeking group'' (Butler, 1997) of middle- and upper-income in-movers. The authors also draw on David Harvey's (1989) account of the shift from urban \emph{managerialism} to \emph{entrepreneurialism} in local governance, and on the urban-entrepreneurial/urban-boosterism literature (Florida's, 2002, ``creative class'' thesis; Swyngedouw et al, 2002; Fainstein, 2008) to frame gentrification as a growth-oriented instrument of interurban competition rather than a housing-provision or welfare policy. Comparative mega-project governance work (Fainstein, 2008; Tasan-Kok, 2010, on the Kop van Zuid specifically) supplies the paper's institutional-comparative lens. The stated theoretical contribution is to shift the analytical focus from \emph{how} such projects are governed to \emph{why} municipal actors pursue gentrification as a goal in the first place -- an angle the authors argue is underdeveloped in the literature.

\subsection*{Data}

\begin{itemize}
  \item \textbf{Source(s)}: Thirty-one semi-structured elite interviews with the principal stakeholders of the two projects (municipal officials, e.g., Rotterdam Development Corporation [OBR] and Department of Housing and Urban Development [DS+V]; Glasgow City Council planners, housing officers and local councillors; private developers, builders, urban designers; representatives of local housing associations), conducted in the interviewees' mother tongues, averaging seventy minutes; supplemented by documentary analysis of municipal policy texts (Rotterdam's \emph{Stadsvisie}/Urban Vision, 2007; \emph{Kop van Zuid} master plans, 1995 and 1999; Glasgow's City Plan and City Plan 2, 2008; Glasgow Clyde Waterfront Regeneration Annual Report, 2005).
  \item \textbf{Geographic scope}: Rotterdam, the Netherlands (Kop van Zuid, a mixed-use waterfront development on the River Nieuwe Maas opposite the city centre) and Glasgow, Scotland, UK (Glasgow Harbour, a brownfield waterfront site on the north bank of the Clyde, roughly 3 km west of the city centre, adjacent to the Partick neighborhood).
  \item \textbf{Spatial unit}: Two flagship neighborhood/project-level waterfront redevelopment sites, nested within a city-level comparative case design contrasting the Dutch and UK/Scottish planning-institutional systems.
  \item \textbf{Time period}: Interviews conducted January 2008--February 2009 (Rotterdam) and September 2008--January 2009 (Glasgow); the projects' planning histories are traced from the late 1980s (Kop van Zuid) and the 1990s (Glasgow Harbour) through to 2011.
  \item \textbf{Sample size}: 31 semi-structured interviews across the two case studies (not disaggregated by city in the text), plus an unspecified number of municipal and project-level policy documents.
  \item \textbf{Key variables}: There are no quantitative dependent/independent variables; the analytical constructs are qualitative -- the presence/explicitness of a municipal ``goal of gentrification,'' the composition and hierarchy of actors (public-sector-led versus developer-led), the housing-stock differentiation strategy (unit type, tenure, price segment), and the framing of gentrification within municipal competitiveness discourse (e.g., ``footloose residents'').
  \item \textbf{Data limitations}: Not reported by the authors as such, but the sample of interviewees is composed entirely of project insiders and policy/development elites; no residents of the new-build developments, displaced or non-displaced residents of adjacent neighborhoods, tenant advocates, or opposition groups were interviewed. Policy documents analyzed are produced by the same municipal and development actors whose rationales are under study.
\end{itemize}

\subsection*{Methodology}

\begin{itemize}
  \item \textbf{Empirical strategy}: Cross-national comparative case-study design (most-similar-systems logic: two former industrial/maritime port cities with parallel economic trajectories of growth, deindustrialization and post-2000 waterfront-led revival), combining semi-structured elite interviews with qualitative documentary/policy-discourse analysis. This is a qualitative, interpretive research design; it is not a causal-inference or econometric study, and the paper reports no regression coefficients, hedonic estimates, or spatial-econometric models.
  \item \textbf{Identification}: Not applicable in the causal-inference sense. The "identifying" logic is comparative-institutional: holding constant broadly similar economic histories and waterfront-redevelopment contexts across the two cities in order to isolate the effect of differing political-economic/planning-institutional arrangements (Dutch top-down, municipally led planning versus the UK's developer-led model with municipal planning support only) on how the shared municipal goal of gentrification is formulated and realized.
  \item \textbf{Spatial treatment}: Spatial dependence is not modeled statistically. Spatial/scalar reasoning is qualitative and institutional: the paper situates each project within its city-level governance context and within a broader national planning-system comparison (Dutch strong master-planning versus Scottish/UK weaker, developer-driven master plans), consistent with the geography domain's emphasis on place specificity and multi-scalar context (Harris, 2008).
  \item \textbf{Robustness checks}: Data triangulation between interview accounts and municipal/project policy documents; use of a two-case comparative design intended to reveal how a shared global strategy (new-build gentrification) is locally inflected by different actor compositions and institutional hierarchies.
  \item \textbf{Methodological innovation}: The paper's distinctive move is not a new statistical or spatial method but an analytical one: examining new-build gentrification through the self-reported rationale of the policymakers and developers who pursue it, and through their internal organizational logic (rendered visually as process diagrams: a top-down municipal model for Rotterdam, Figure 1, versus a developer-led, planning-support model for Glasgow, Figure 2), rather than through governance-structure analysis alone (cf. Tasan-Kok, 2010) or through outcome-side housing-market indicators.
\end{itemize}

\subsection*{Key Findings}

\begin{enumerate}
  \item \textbf{Explicit versus implicit municipal goal-setting}: Rotterdam's municipality names gentrification as an explicit policy objective in its 2007 \emph{Stadsvisie} (Urban Vision), openly framing the retention of ``footloose'' affluent and middle-class residents as necessary to counteract the perceived liability of the city's historically high share of social-rented housing. Glasgow's equivalent goal is present but far more implicit and blurred with an undifferentiated aim of attracting \emph{any} population growth (including asylum seekers under the Home Office dispersal programme), rather than a targeted higher-income cohort. Magnitude/significance: qualitative/discursive finding, not statistically quantified; supported by direct textual and interview evidence (e.g., the Urban Vision passage on ``a balanced composition of the population,'' Gemeente Rotterdam, 2007, page 63).
  \item \textbf{Divergent governance pathways, convergent outcome}: The Dutch, municipally led planning model (Kop van Zuid: city sets the master plan; private developers execute a municipal mandate) and the UK developer-led model (Glasgow Harbour: Clydeport/Glasgow Harbour Ltd drives the project; the council's role is limited to planning support and light-touch facilitation) both culminate in state-sanctioned new-build gentrified waterfront housing, but via different causal chains -- ``state-led'' in Rotterdam versus ``state-sponsored, developer-led'' in Glasgow (the authors' own Figures 1 and 2). Magnitude/significance: process-tracing finding based on interview and documentary evidence, not a statistical estimate.
  \item \textbf{Shared rationale: housing-stock mismatch and interurban competitiveness}: In both cities policymakers justify the strategy by pointing to a perceived shortage of higher-income/owner-occupied housing stock relative to competing municipalities, and frame the retention of affluent, skilled ("footloose") residents as essential to urban economic competitiveness and tax-base/spending-power gains (e.g., the Glasgow planner's comment that incoming residents mean "more spending power," "more people employed in the retail and restaurant industries"). The evidence supports this as the authors' central, well-triangulated interpretive claim across both case studies.
  \item \textbf{Rotterdam judged more successful in realizing the goal}: The authors argue -- based on the greater diversity of housing types delivered (family homes, apartments, single-family dwellings with gardens) -- that gentrification was achieved more completely and more diversely in the Kop van Zuid than in Glasgow Harbour, where the housing product (predominantly one- and two-bedroom flats) and the resulting enclave-like, poorly city-integrated development reflect weaker municipal goal-setting and looser master planning. This is presented as the authors' interpretive assessment rather than a measured outcome (e.g., no post-occupancy demographic or price data are reported).
  \item \textbf{Denial of displacement in official discourse}: Stakeholders in both cities, and especially in Glasgow, deny that new-build gentrification produces displacement, on the grounds that the sites were vacant industrial/brownfield land rather than occupied residential neighborhoods (illustrated by the local councillor's remark that "there wasn't the case of a community being displaced''). The authors treat this claim as a self-serving element of an emerging, depoliticized "positive" gentrification discourse (cf. Wyly and Hammel, 2008) rather than as an independently verified absence of exclusionary or indirect displacement.
\end{enumerate}

\subsection*{Contributions}

\begin{enumerate}
  \item \textbf{Empirical contribution}: Provides one of the first systematic, interview-based comparative accounts of \emph{why} municipal governments in two second-tier, post-industrial European port cities actively pursue new-build gentrification, documented directly through the discourse of the policymakers, planners, and developers responsible for it.
  \item \textbf{Theoretical contribution}: Extends the "new-build gentrification" concept (Davidson and Lees, 2005) and Harvey's (1989) managerialism-to-entrepreneurialism framework by showing that the same global strategy of state-supported new-build gentrification can be realized through structurally opposite actor hierarchies (state-led versus developer-led) while converging on a shared underlying goal -- countering a city's self-perceived status as ``a rich city with poor people.''
  \item \textbf{Methodological contribution}: A cross-national, most-similar-systems comparative case-study design built on triangulated elite interviews and policy-document analysis, applied specifically to municipal \emph{rationale} rather than governance structure alone (distinguishing it from Tasan-Kok's, 2010, governance-focused comparison of the same Rotterdam case).
  \item \textbf{Policy contribution}: Cautions that new-build gentrification, however effective at attracting affluent in-movers, is a wealth-\emph{creation} strategy rather than a wealth-\emph{distribution} strategy: it does little to address unemployment, affordable-housing shortages, or the "more pressing urban economic and social problems" of the cities pursuing it, and risks entrenching rather than reducing pre-existing socio-spatial divisions.
\end{enumerate}

\subsection*{Limitations and Critical Assessment}

\begin{enumerate}
  \item \textbf{Identification threats}: This is not a causal-inference study, so identification in the econometric sense does not apply; the relevant threat is interpretive validity. The interview sample is composed almost entirely of project insiders (municipal officials, developers, planners) with a direct professional or institutional stake in presenting the projects favorably. Absent are residents of the new developments, displaced or non-displaced residents of adjacent working-class areas (e.g., Partick, next to Glasgow Harbour), tenant organizations, or critical academics/journalists whom the authors themselves cite as holding opposing views (e.g., Boddy, 2007; McIntyre, 2008; media critics referenced in section 5.2). The near-total reliance on elite, self-interested informants risks systematically reproducing an official, legitimizing narrative rather than independently verifying it -- most visibly in the finding that "there was little official opposition" and the flat denial of displacement, both of which come directly from the stakeholders whose own decisions are being evaluated.
  \item \textbf{External validity}: The study rests on exactly two cases, both mid-sized, historically maritime, Northern European port cities with broadly similar deindustrialization trajectories. The authors' own framing situates the cases within a wider claim about "cities further down the urban hierarchy" globally, but the two-case design cannot support that broader generalization; the findings may not extend to cities with different housing-tenure structures (e.g., predominantly private-rental or homeownership-dominant systems outside northwest Europe), weaker or absent land-use planning authority, or non-Anglo/non-Dutch legal-institutional traditions.
  \item \textbf{Omitted mechanisms}: The paper studies policy \emph{discourse} and \emph{process}, not \emph{outcomes}. It does not report census-tract-level or neighborhood-level demographic, income, or house-price data to verify whether the intended socioeconomic upgrading actually occurred, at what rate, or how it compares to the counterfactual trajectory of the surrounding city. It also does not examine possible spillover or filtering effects on adjacent working-class neighborhoods (e.g., whether new-build gentrification in Glasgow Harbour affects rents or turnover in neighboring Partick), nor the opportunity cost of channeling public planning capacity and (in Rotterdam) direct municipal investment into flagship affluent housing rather than upgrading the existing social-housing stock.
  \item \textbf{Data constraints}: Stronger data would include household-level survey or migration data on who actually moved into the Kop van Zuid and Glasgow Harbour (income, prior residence, tenure history) to test the "footloose resident retention" mechanism directly, rather than inferring it from policymaker intent; longitudinal house-price or rent data for the two sites and their surrounding tracts; and interviews with residents of both developments and adjacent areas to assess indirect/exclusionary displacement empirically rather than relying on stakeholder assertions that none occurred.
  \item \textbf{Equilibrium considerations}: The paper does not address general-equilibrium or metropolitan-wide adjustment effects: whether attracting affluent households to two flagship waterfront enclaves merely redistributes a fixed regional pool of higher-income households (a zero-sum reshuffling within the city-region) rather than generating net inward migration or net new tax base, and whether the resulting bifurcated housing supply (affluent waterfront enclave versus deteriorating existing social-housing stock) is compatible with, or in tension with, each city's broader affordable-housing strategy. The interviews were also conducted in 2008--2009, coincident with the onset of the global financial crisis, yet the paper does not discuss how a housing-market downturn might affect the viability of a developer-led, market-exposed project such as Glasgow Harbour going forward.
\end{enumerate}

\subsection*{Cross-Disciplinary Bridges}

\begin{itemize}
  \item \textbf{Connection to Urban Economics}: The municipal narrative of a housing-stock "mismatch" (too little owner-occupied, higher-quality housing relative to competitor suburban municipalities) echoes urban-economics accounts of housing supply constraints and filtering (Glaeser and Gyourko; Saiz, 2010), but the paper treats this claim as discourse to be documented and critically examined rather than as a quantity to be estimated with supply-elasticity or hedonic methods.
  \item \textbf{Connection to Geographic Finance}: Glasgow Harbour's developer-led model, in which a former port-management subsidiary (Clydeport/Glasgow Harbour Ltd) converts brownfield land into a profit-oriented, market-rate residential asset with no social-housing provision, is a clear instance of the financialization-of-housing logic described by Aalbers (2016) -- housing production driven primarily by capital-accumulation motives rather than by a public housing mandate, with the municipality's role reduced to a facilitator of private investment.
  \item \textbf{Connection to Regional Science}: The competitiveness framing used by both cities' policymakers -- attracting and retaining skilled, "footloose" residents to sustain the local tax base and service economy -- parallels the logic of Rosen-Roback/Moretti spatial-equilibrium and human-capital-attraction models and the place-based-policy debate (Kline and Moretti, 2014), though tested here through qualitative discourse analysis rather than formal spatial-equilibrium estimation.
  \item \textbf{Suggested related works}:
    \begin{itemize}
      \item Davidson, M. and Lees, L. (2005) -- "New-build `gentrification' and London's riverside renaissance," the direct conceptual antecedent for treating flagship new-build housing as a mutation of gentrification.
      \item Smith, N. (1979) -- "Toward a theory of gentrification: a back to the city movement by capital, not people," the foundational rent-gap framework invoked to justify new-build investment in "blighted" waterfront land.
      \item Tasan-Kok, T. (2010) -- "Entrepreneurial governance: challenges of large-scale property-led urban regeneration projects," a governance-focused comparative study of the same Kop van Zuid case that this paper's actor-rationale analysis complements.
    \end{itemize}
\end{itemize}

\subsection*{Quotable Passages}

\begin{quote}
"The problem with Rotterdam is that it's a rich city with poor people." (p.~1438)
\end{quote}

\begin{quote}
"[Y]ou have to have a good job if you want to live here.\ldots{} People without a job, or with a bad job, or a relatively low salary, they cannot live on the Wilhelmina Pier, no. But they also can't live in Hollywood." (p.~1447)
\end{quote}

\begin{quote}
"[W]hat's wrong is if at the end it denies other people access to housing, but that actually hasn't been, it isn't the case with the Harbour development. Because it wasn't the case of a community being displaced so that these flats could be built, it was urban industrial land." (p.~1449)
\end{quote}
